% !PRINTED 423
\weberpage{423}

%<<CONT
einführen, so ist (\S.\,46)
\[
D = P^2.
\]
Es ist also auch nach (8)
\weq{10}{P = n^{\frac{n-3}{2}} \sqrt{(-1)^{\frac{n-1}{2}} n},}
wodurch $P$, abgesehen vom Vorzeichen, bestimmt ist. Es ist $P$ reell, wenn $n \equiv 1$, und imaginär, wenn $n \equiv -1 \ (\mathrm{mod}\ 4)$ ist, also reell für $n = 5, 13, 17, 29, 37 \ldots$, imaginär für $n = 3, 7, 11, 19, 23, 31 \ldots$ Das Vorzeichen, das wir in (10) der Wurzel zu geben haben, hängt davon ab, welches $r$ wir in dem Ausdruck (9) gewählt haben.

Wählen wir ein bestimmtes $r$, z. B. $r = r_0$, so können wir das Vorzeichen in (10) noch bestimmen. Um dies auszuführen, theilen wir die binomischen Factoren von $P$,
\[
r^\nu - r^\mu, \quad \nu < \mu,
\]
deren Anzahl $\frac{n(n-1)}{2}$ beträgt, in zwei Classen. Die Differenzen der einen Classe bilden das Product
\weq{11}{Q = (r - r^{n-1})(r^2 - r^{n-2}) \cdots \left(r^{\frac{n-1}{2}} - r^{\frac{n+1}{2}}\right),}
das also alle die Factoren enthält, in denen $\mu + \nu = n$ ist.

Die übrigen Factoren lassen sich in Paaren von folgender Form zusammenfassen
\weq{12}{R = (r^\nu - r^\mu)(r^{n-\mu} - r^{n-\nu}).}
Die Anzahl der Factoren in $Q$ ist $\frac{n-1}{2}$, und also ist die Anzahl der Paare $R$
\[
\frac{1}{2}\left(\frac{n(n-1)}{2} - \frac{n-1}{2}\right) = \left(\frac{n-1}{2}\right)^2.
\]
Nun ist, wenn
\[
r = e^{\frac{2\pi i k}{n}}
\]
ist,
\[
R = -2 + r^{\mu-\nu} + r^{\nu-\mu} = -2\left(1 - \cos\frac{2\pi(\mu-\nu)}{n}\,k\right),
\]
d. h. $R$ ist immer negativ. Folglich hat das Product aller $R$ das Vorzeichen
\[
(-1)^{\frac{n-1}{2}}.
\]
Es ist ferner
\weq{13}{Q = (2i)^{\frac{n-1}{2}} \sin\frac{2\pi k}{n} \sin\frac{4\pi k}{n} \cdots \sin\frac{(n-1)\pi k}{n},}
%CONT>>
